\newpage
\section{Discussion}
\label{sec:discussion}

% Chapter purpose: Interpret the structural ceiling alongside quality
% evaluation dependencies and module-space limitations.  Lead with
% information asymmetry, then lambda dominance, then HyDE resolution,
% then quality evaluation, then methodology.  State limitations and
% future directions.


\subsection{Summary of Findings}
\label{subsec:synthesis}

This thesis investigated whether per-query module activation can reduce
energy consumption in locally deployed RAG systems without degrading answer
quality.  Three results, each building on the previous, converge on a
single conclusion.

The baseline analysis (Chapter~\ref{sec:baseline_analysis}) confirmed that
always-on module activation is measurably inefficient.  The full pipeline
costs 5.6$\times$ more energy than the minimal pipeline, while
quality improves by less than 0.06~points on either metric.  The cost is unconditional;
the benefit is conditional and query-dependent.

The decision problem formulation (Chapter~\ref{sec:controller}) quantified
the opportunity.  In the full eight-configuration space, an oracle router
saves 75\% energy versus always-C7(KMS Pipeline) by avoiding HyDE whenever possible.
Restricting to the non-HyDE configurations actually tested for routing
(\{C0, C2, C3, C6\}), the oracle saves 19\% versus always-C6 and 40\%
of queries benefit from routing to a cheaper configuration.  Yet
pre-execution query features explain near-zero variance in quality scores
($R^2 \approx 0$), signalling that even this reduced opportunity may be
difficult to exploit in practice.

The routing investigation (Chapter~\ref{sec:strategies}) confirmed this
difficulty.  Five routing strategies and three explorations, spanning
three model classes, two feature sets and two decision
architectures, all converge to a config accuracy of
0.38--0.40.  No approach produces a statistically meaningful improvement
over the cheapest-configuration baseline (Static~C0 at 0.382).  The
ceiling is signal-limited, not model-limited.


\subsection{Information Asymmetry}
\label{subsec:ceiling-interpretation}

The convergence of all routing strategies to the same narrow performance
band is the thesis's central empirical finding.  The 60.5~percentage-point
gap between oracle (100\%) and best router (39.5\%) is not a failure of
modeling.  It quantifies information that exists at retrieval time but is
invisible before pipeline execution.  This section interprets the nature
of this asymmetry and what it implies for the field.

\paragraph{A circular dependency.}
Accurate routing requires post-retrieval signals: document scores,
score gaps, source diversity.  The purpose of routing is to avoid
executing retrieval variants that would produce those signals.  This
creates a circular dependency: the information needed to route is
produced by the very computation that routing aims to skip.

Experiment E3 (Section~\ref{subsec:explorations})
tested this directly.  Even after observing C0 retrieval statistics,
the Stage~2 reranking classifier achieved AUC~$\approx 0.53$, barely
above random.  The retrieval statistics answer ``how good was this
retrieval?'' but not ``would a different retrieval strategy help this
query?''  These are structurally different questions.

\paragraph{Module-activation vs.\ model routing.}
Table~\ref{tab:routing-comparison} contrasts three routing granularities.
Model routing selects \emph{who answers}: the query goes to a strong or
weak LLM with the same context.  Module-activation routing selects
\emph{how to prepare the context}: which retrieval and post-retrieval
stages run before generation.  The quality gap that a router must
discriminate shrinks by an order of magnitude from model selection
(10--30~pp) to module activation (2--4~pp), which directly limits the
signal available to any classifier.

\begin{table}[h!]
\centering
\small
\begin{tabular}{l l c l}
\hline
\textbf{System} & \textbf{Granularity} & \textbf{Quality gap} & \textbf{Reported savings} \\
\hline
RouteLLM~\cite{routellm2025}   & Model selection   & 10--30~pp & 50--85\% cost reduction \\
FrugalGPT~\cite{frugalgpt2023} & Model cascade     & 10--30~pp & up to 98\% cost reduction \\
This work                       & Module activation & 2--4~pp   & 39.5\% config accuracy \\
\hline
\end{tabular}
\caption{Routing approaches by granularity.}
\label{tab:routing-comparison}
\end{table}

In model routing, a router discriminates between generators that produce
10--30~pp quality gaps on identical context.  In our module routing setup,
the same generator receives slightly different retrieval preparations,
compressing the quality gap to 2--4~pp while within-config noise stays
constant.  Systems with larger module-level quality variation may face a
more tractable routing problem.  Adaptive RAG systems
(Self-RAG~\cite{selfrag2024}, CRAG~\cite{crag2024},
Adaptive-RAG~\cite{adaptiverag2024}) operate at the strategy level where
quality differences remain large enough to learn.  In our test case,
finer-grained module control faces a structurally harder prediction
problem.

\paragraph{Signal-limited, not model-limited.}
Because the ceiling is signal-limited, investing in more sophisticated
classifiers (deeper networks, larger training sets, richer architectures)
is unlikely to yield improvements within the current feature
space.  The between-experiment accuracy spread (2~pp) is comparable to
within-experiment cross-validation noise (1--2~pp).  The rational response
is to deploy the simplest approach that captures the available signal and
invest research effort in cheaper ways to access post-retrieval
information.

\paragraph{Characterising the boundary.}
Defining a ceiling has methodological value.  The experiments
consistently place the boundary at 0.38--0.40 config accuracy for
pre-execution features on module-activation routing, empirically
validated across five strategies and three explorations.  This delimits
where future work is unlikely to yield gains under similar assumptions
and identifies what is needed to break through: cheaper access to
post-retrieval signals.  The 60.5~pp gap between oracle and achieved
performance defines a concrete, quantified research target.


\subsection{Energy Penalty ($\lambda$) Analysis}
\label{subsec:lambda}

Perhaps the most actionable finding is that the energy penalty
parameter~$\lambda$ matters more than the choice of classifier, feature
set or model architecture.  Across all experiments, the best results
appear at $\lambda = 1.25$--$1.50$.  Without the energy penalty
($\lambda = 0$), all models achieve 0.28--0.35, worse than Static~C0.
Performance rises only as $\lambda$ increases, which pushes selections
toward cheap configurations unless the classifier is confident.  In
practice, this amounts to a soft version of ``always pick C0'': the
utility function, not the classifier, drives accuracy above baseline.

This observation shifts the RAG efficiency problem from a machine learning
problem to a policy problem.  A high~$\lambda$ implements a
cheap-by-default policy: expensive modules are assumed unnecessary unless
the quality signal strongly indicates otherwise.  For locally deployed
systems with constrained energy budgets, this policy is the most
defensible strategy.  It requires no additional model, no training
infrastructure and no runtime overhead beyond the utility calculation.

The classifier's marginal contribution is the 1--2~pp it adds on top of
the $\lambda$-driven default.  Tuning a single scalar parameter
outperforms all model-selection experiments in this thesis.


\subsection{The HyDE Variance}
\label{subsec:hyde}

Hypothetical Document Embeddings (HyDE)~\cite{hyde2023} is widely
reported as beneficial in RAG pipelines, particularly for bridging
vocabulary gaps between queries and documents.  Our results contradict
this: HyDE degrades quality ($-$0.013~fc, $-$0.021~faithfulness) at
5.3$\times$ the energy cost of the minimal pipeline.  This led to its
exclusion from the routing configuration space.

The contradiction likely reflects an interaction between HyDE and the
retrieval architecture.  HyDE was designed for sparse retrieval systems
where lexical mismatch dominates.  In our dense-retrieval setup using
BGE-M3, the embedding model already bridges vocabulary gaps.  We
hypothesise that the hypothetical document generated by the LLM pushes
the query embedding into semantically adjacent but factually irrelevant
regions of the vector space.  Where sparse retrieval benefits from
vocabulary expansion, dense retrieval suffers from what we term
\emph{semantic drift}: the hypothetical introduces noise into an already
well-positioned query vector.  Confirming this hypothesis requires
comparing query and hypothetical-document embeddings directly in the
vector space and remains an open direction for future work.

This result does not invalidate HyDE as a technique.  It demonstrates
that module utility is system-dependent: a component that improves one
retrieval architecture may degrade another.  The 75\% oracle energy
saving in the full configuration space is primarily an engineering
decision (disabling a harmful module), while the 19\% saving in the
non-HyDE routing space is the algorithmic ceiling for per-query routing.
These two numbers quantify different opportunities: the first is
immediately actionable, the second defines the research target.


\subsection{Methodological Legacy}
\label{subsec:methodology}

The measurement framework developed in
Chapter~\ref{sec:method_protocol} is the most transferable artifact of
this thesis.  The framework implements three design principles:
session-scoped energy counters, pipeline-stage markers and per-node
attribution.  These generalise to any distributed system with toggleable
components.  The instrumentation code, configuration API and analysis
pipeline are available as part of the CITI~KMS repository.  While the
absolute energy numbers are system-specific, the methodology enables any
modular RAG deployment to quantify its own trade-offs and assess routing
potential without modifying pipeline stages.

\paragraph{Quality evaluation.}
The routing pipeline depends on quality labels, which depend on evaluation
metrics, which depend on the evaluator model.  This dependency chain
(evaluator $\to$ thresholds $\to$ labels $\to$ routing signal) means
that evaluation quality propagates through every downstream result.
If the labels are unreliable, no classifier can learn a meaningful
decision boundary regardless of feature strength.

We use Qwen2.5-14B-Instruct as evaluator and a different model family
(R-4B) as generator, avoiding self-evaluation bias where models favour
their own generation patterns~\cite{vrageval2024}.  Wang~et~al.\ report
that evaluator size determines reliability: 8B~models achieve 36.8\%
human agreement on binary accept/reject versus 82.6\% for
GPT-4-class models.  Our 14B~evaluator sits between these extremes.
The choice of RAGAS was driven by its existing integration within the
system infrastructure; methodologically, it operationalises the
LLM-as-a-judge paradigm by decomposing answers into verifiable
claims~\cite{ragevalsurvey2025}.  Any inconsistency in the judge's
scoring propagates directly into the routing labels.  A comparison with
alternative frameworks (ARES, vRAG-Eval) would verify whether the
structural ceiling is judge-invariant.

The ICC analysis in Section~\ref{subsec:quality-labeling} shows that each
metric privileges a different aspect of the pipeline.  Answer~F1
(ICC\,=\,0.825) captures lexical overlap but treats all configurations as
near-interchangeable, making it unsuitable for routing supervision.
Faithfulness (ICC\,=\,0.432) has the strongest between-configuration
variance but evaluates generation consistency rather than factual
accuracy.  Factual correctness (ICC\,=\,0.583) balances the two but
depends on the evaluator model's decomposition of claims.  The E2
exploration (Section~\ref{subsec:explorations}) confirms that the metric
choice is decisive: only factual correctness combined with faithfulness
at $\tau = 0.7$ produces a positive routing margin; all other combinations
collapse to baseline or below.

For a knowledge management deployment, the priority is whether the answer
contains the correct information (factual correctness) and does not
fabricate unsupported claims (faithfulness).  Lexical overlap with a
reference answer matters less when users seek explanations rather than
exact phrases.  The structural ceiling of 0.38--0.40 config accuracy may
therefore reflect not only information asymmetry but also label noise from
imperfect evaluation.  These two sources of difficulty compound: even if
corpus-aware features provided stronger signal, noisy labels would limit
what a classifier can learn from them.  Establishing a validated,
domain-appropriate quality evaluation method is the logical first step
before investing in richer routing features.


\subsection{Limitations}
\label{subsec:limitations}

Several limitations constrain the generalisability of these findings.

\paragraph{System generalisability.}
All experiments use the CITI~KMS pipeline with a specific LLM size,
embedding model, retrieval backend and three toggleable modules (HyDE,
reranking, hybrid search) inherited from the existing system rather than
selected from a systematic survey.  Advanced RAG architectures offer
additional modules---query decomposition, web-search fallback
(CRAG~\cite{crag2024}), self-reflective generation
(Self-RAG~\cite{selfrag2024}), document summarisation, answer
verification---each with different energy costs.  The 2--4~pp quality gap
may be specific to this module set; a broader module space could produce
a more tractable routing problem.  All measurements also assume
single-request sequential execution; concurrent workloads may change
energy patterns due to resource contention and GPU utilisation differences.
The absolute numbers (e.g., 5.6$\times$ energy ratio) are
system-specific; the methodology and ceiling characterisation approach
generalise.

\paragraph{Dataset scope and feature transferability.}
HotpotQA provides multi-hop reasoning queries only, and the experiment
uses a corpus built from gold supporting documents.  In open-domain
retrieval, where relevant documents may be absent or diluted, the quality
surface may differ and could create stronger routing signals.  The two
most important routing predictors are dataset-specific:
\texttt{q\_has\_logical\_operator} appears in 61\% of HotpotQA queries
versus 22\% of KILT queries (39~pp gap), and \texttt{q\_multi\_entity}
in 76\% versus 45\% (31~pp gap).  Both features reflect multi-hop
structure rather than a general query property.  A router trained on
these features would encounter a substantially different input
distribution on factoid or verification queries, limiting
transferability without retraining.

\paragraph{Title-level provenance evaluation.}
HotpotQA annotates supporting facts as (title, sentence\_id) pairs; the
official Sup~F1 operates at sentence granularity.  Our system chunks
Wikipedia articles by character count (1{'}024-character windows with
200-character overlap), not by sentence boundaries.  Our Sup~F1 evaluates
at title level and is not directly comparable to leaderboard scores.

\paragraph{LLM-as-judge evaluation.}
Quality evaluation relies on RAGAS with Qwen2.5-14B-Instruct as the
evaluator model.  LLM-based evaluation introduces its own biases.  We
mitigate this by using the same evaluator consistently across all
configurations, so relative comparisons remain valid even if absolute
scores carry systematic bias.

\paragraph{Limited feature set.}
The 23~pre-execution features are a first attempt at characterising query
difficulty for routing.  Richer features (query--document interaction
signals, LLM self-assessment or lightweight retrieval probes) may carry
stronger discriminative power, though the ceiling analysis suggests this
is unlikely without accessing post-retrieval information.


\subsection{Future Directions}
\label{subsec:future}

The ceiling finding defines the research agenda: how to access
post-retrieval information without incurring its full cost.

\paragraph{Lightweight retrieval probing.}
The circular dependency can be partially broken if a cheap retrieval
approximation provides routing-relevant signal.  A sparse retrieval
probe (e.g., BM25) that costs milliseconds rather than seconds could
estimate whether hybrid search or reranking would help a specific query.
If a BM25 score gap predicts reranking benefit, the router gains access
to post-retrieval signals at a fraction of the pipeline cost.

\paragraph{Transfer to model routing.}
Module-activation routing faces quality gaps of 2--4~percentage points;
model routing faces gaps of 10--30~pp.  Investigating whether the same
query features that fail for module routing succeed for model routing
would clarify the relationship between signal strength and routing
granularity.  Success would confirm that the ceiling is specific to
fine-grained module control, not to query-feature-based routing in
general.

\paragraph{Multi-task workloads.} A production RAG system serves diverse query types.  Different query
types (factoid, verification, multi-hop) may create more separable
routing populations.  Routing accuracy may improve when the training
set spans multiple task distributions with larger between-type
quality variation.

\paragraph{Non-routable query detection.} The 445~queries (37\%) that fail under all eight configurations represent
a system quality limitation, not a routing limitation.  If non-routable
queries are detectable pre-execution, a router could abstain rather than
waste energy on queries likely to fail regardless.  This binary
distinction may be more tractable than configuration selection, since the
routable/non-routable gap is larger than the 2--4~pp quality gap between
configurations.

\paragraph{Online learning.} Adapting the routing policy from production feedback using contextual
bandits or reinforcement learning would enable continuous improvement
without requiring forced-run experiments.  This is particularly relevant
because the quality landscape may shift as the document corpus evolves.

\paragraph{Corpus-aware embedding features.}
The 23~pre-execution features characterise queries in isolation from the
indexed document collection.  Since BGE-M3 embeddings already exist for
all indexed documents, principal component analysis of the corpus
embedding space could yield routing features at negligible cost.
Projecting each query into the first 100--150 principal components of the
document cloud enables features such as distance to the corpus centroid,
local document density and cluster membership.  UniRoute~\cite{uniroute2025}
provides theoretical grounding for cluster-based routing, demonstrating
Bayes-optimal routing with bounded excess risk when queries are assigned
to corpus clusters.  These features partially break the circular
dependency (Section~\ref{subsec:ceiling-interpretation}) by providing
corpus awareness without executing retrieval variants.  They remain
pre-execution, so they cannot access actual retrieved document quality,
but they represent a middle ground between query-only surface features
and the full cost of post-retrieval signals.

\paragraph{Broader module evaluation.}
The three modules tested in this work span query expansion, retrieval
strategy and post-retrieval filtering.  Production RAG systems may
include additional stages: query decomposition for complex questions,
web-search fallback for corpus gaps, answer verification against retrieved
evidence or document summarisation before generation.  Each module
introduces a different energy--quality trade-off.  Systematically
evaluating the quality contribution of each candidate module---before
attempting to route among them---would establish which modules produce
quality differences large enough to be learnable by a pre-execution
router.

\paragraph{Energy measurement under concurrent workloads.}
Extending the measurement framework to multi-tenant scenarios would
quantify how resource contention affects per-request energy attribution
and whether routing decisions remain valid under production load.
